% Chapter 3, section 2: Publish/subscribe and fan-out
% Write in plain LaTeX. Code of a guide step: \guide{<step>}; commands: \begin{shell}...\end{shell};
% a screenshot: \screenshot{file}{caption}; a gap in the draft: \placeholder{...}.
% See book/latex/queen.sty for the full list.
%
% Draft by Claude, 2026-10-24. The code is cut from examples/tutorials/js/book-use-cases.mjs, which
% checks every result quoted here against a 2.7.0 node.

\section{Publish/subscribe and fan-out}

One event, many readers: when an order is placed, billing charges it, shipping prepares it, and
analytics counts it. Each reader must get every event, at its own pace, and a reader added later
must be able to start from the beginning.

Each reader is a consumer group. A group gets every message of the queue and keeps its own
position; the workers inside a group share its work, as in a work queue.

\queenfig{pubsub}{Published once, read by every group, also by one that joins later.}

\guide{uc-pubsub}

Billing and shipping each got the three orders. Analytics, started later, read all three from the
start.

\why{the queue is a log: reading does not delete. So a new subscriber costs nothing to the
publisher, and a group that falls behind slows down only itself.}

One thing to know: a new group starts at the end of the queue unless it asks for
\texttt{subscriptionMode('all')}, or a point in time with \texttt{subscriptionFrom}. Our check
proved it: a group created without it found nothing.
